%%%%%%%%%%%%%%%%%%%%%%%%%%%%%%%%%%%%%%%%%
% Academic Paper Response Letter. 
% LaTeX Template
% Version 1.0 (20/07/2023)
% 
% This template refers to the toolbox tutorial (in Chinese):
% https://liam.page/2016/07/22/using-the-tcolorbox-package-to-create-a-new-theorem-environment/
% Please refer to the tutorial for your customization. 
% 
% Template license: Creative Commons CC BY 4.0
%
% Author: Yuhan Xie
% Contact email: yhxie0107@gmail.com
% Any revision advice is welcome!
% 
%%%%%%%%%%%%%%%%%%%%%%%%%%%%%%%%%%%%%%%%%


\documentclass{article}
%% Language and font encodings
\usepackage[english]{babel}
\usepackage{times} % make the section title in Times New Roman

\usepackage{booktabs}
\usepackage{tabu}
\usepackage[T1]{fontenc}

%% Sets page size and margins
\usepackage[a4paper,top=2.5cm,bottom=2cm,left=1.7cm,right=1.7cm,marginparwidth=1.75cm]{geometry}

%% Useful packages
\usepackage{amsmath}
\usepackage{graphicx}
\usepackage[colorinlistoftodos]{todonotes}
\usepackage[colorlinks=true, allcolors=blue]{hyperref}
\usepackage{indentfirst}
\usepackage{graphicx}
\usepackage{subfigure}
\usepackage{float}
\usepackage{threeparttable} 
\usepackage{multirow}
\usepackage{url}


\usepackage[most]{tcolorbox}


%% Comment box environment used throughout the letter
\newtcbtheorem[auto counter]{cmt}{Comment}{
	colbacktitle = black!60!white, colframe = black!60!white,
	colback = black!5!white,
	fonttitle=\bfseries,%fontupper=\itshape,
}{t}


\begin{document}

%%%%%%%%%%%%%%%%%%%%%%%%%%%%%%%%%%%%%%%%%%%%%%%%%%%%%%%%%%%%%%%%%%%%%%%%%%%%%%%
% Response to Reviewers -- PONE-D-26-28965
%
% Structure follows POINT_BREAKDOWN.md (preliminary-structure version):
% four sections, no separate "Response to the Editor" section --- the
% journal-side items are grouped under "Journal Requirements".
%   1. Journal Requirements ........  7 items
%   2. Reviewer #1 ................. 11 items
%   3. Reviewer #2 ................. 14 items
%   4. Reviewer #3 ................. 12 items
%   Total ......................... 44 response items
%
% Comment boxes quote the review verbatim. The state of each response is marked
% visibly in the body of the letter (see the banner on the first page), never
% here: %-comments do not reach the PDF that the authors actually read.
%%%%%%%%%%%%%%%%%%%%%%%%%%%%%%%%%%%%%%%%%%%%%%%%%%%%%%%%%%%%%%%%%%%%%%%%%%%%%%%


\begin{center}
{\Large\bfseries Response to Reviewers}

\vspace{0.2cm}
Manuscript PONE-D-26-28965

\vspace{0.1cm}
\textit{Evaluation of validation strategies and transferability of EEG attention
decoding: a comparative analysis on open datasets}
\end{center}

\vspace{0.3cm}

\noindent
\textbf{[TODO --- working state of this letter; this banner is not part of the
submission.]} Every response below is in one of three states. A response with no
marker is backed both by the code in \texttt{python\_project/} and by text that is
actually present in \texttt{manuscript\_clean/manuscript.tex}. A response
followed by a \textbf{[TODO]} paragraph is only partly backed: the
implementation exists, but the manuscript does not yet say so, or the numbers it
refers to have not been computed. A response consisting only of \textbf{[TODO]}
is not answered at all yet. The manuscript is mid-rewrite --- its Abstract,
Background, Models and Results sections are cleared placeholders --- the full
sweep in \texttt{python\_project/results/} is still running, and no analysis has
been produced from it. Until that changes, no response here may cite a number, a
table, a figure or a manuscript passage. All markers must be gone before the
letter is uploaded.

\vspace{0.3cm}
\noindent
Dear Editor and Reviewers,

\vspace{0.2cm}
\noindent
We thank the Editor and the Reviewers for their careful evaluation of our
manuscript and for their constructive comments. We have carefully considered all
comments and are preparing a revised version of the manuscript that addresses
the issues raised during peer review.

\vspace{0.2cm}
\noindent
The major revisions will include:

\begin{itemize}
	\item \textbf{[TODO --- major revision 1: redefinition of the target
	      construct.]} The direction is settled --- the study claims
	      discrimination between operationally defined experimental conditions
	      rather than decoding of attention as a psychological construct --- but
	      the manuscript does not yet carry it: Table~1 still names its classes
	      ``High Attention'' and ``Low Attention'', and no label-mapping
	      sensitivity analysis exists in the code. Only the exclusion of drowsy
	      epochs from the primary analysis is in place. Rewrite this bullet
	      together with the responses to MR1 and Comments~3.1, 3.2 and 4.5.
	\item \emph{[Major revision 2 --- to be completed after the corresponding manuscript changes are finalized.]}
	\item \emph{[Major revision 3 --- to be completed after the corresponding manuscript changes are finalized.]}
	\item \emph{[Additional major revision, if needed.]}
\end{itemize}

\noindent
For clarity, our response is organized into four sections:

\begin{enumerate}
	\item \textbf{Journal Requirements}
	\item \textbf{Reviewer \#1}
	\item \textbf{Reviewer \#2}
	\item \textbf{Reviewer \#3}
\end{enumerate}

\noindent
Reviewer and journal comments are reproduced below and followed by our
point-by-point responses. Where a single original paragraph contains several
independent criticisms or requests, it is divided into separate response points
so that each issue can be addressed explicitly. Changes to the manuscript are
identified with the corresponding page and line numbers wherever applicable.

\vspace{0.4cm}
\noindent
Yours sincerely,

\vspace{0.2cm}
\noindent
\emph{[TODO: author list]}


%%%%%%%%%%%%%%%%%%%%%%%%%%%%%%%%%%%%%%%%%%%%%%%%%%%%%%%%%% Journal Requirements
\newpage
\setcounter{tcb@cnt@cmt}{0}
\section*{Journal Requirements}

\noindent
The seven requirements listed in the decision letter are addressed below.

\vspace{0.3cm}

%%%%%%%%%%%%% Journal Requirement 1
\begin{cmt}{}{}

Please ensure that your manuscript meets PLOS ONE's style requirements,
including those for file naming. The PLOS ONE style templates can be found at
\url{https://journals.plos.org/plosone/s/file?id=wjVg/PLOSOne_formatting_sample_main_body.pdf}
and
\url{https://journals.plos.org/plosone/s/file?id=ba62/PLOSOne_formatting_sample_title_authors_affiliations.pdf}

\end{cmt}
\vspace{0.1cm}
\noindent
\underline{\textbf{Response:}}
\vspace{0.2cm}

\emph{[TODO]}

\vspace{0.3cm}


%%%%%%%%%%%%% Journal Requirement 2
\begin{cmt}{}{}

Please note that PLOS One has specific guidelines on code sharing for
submissions in which author-generated code underpins the findings in the
manuscript. In these cases, we expect all author-generated code to be made
available without restrictions upon publication of the work. Please review our
guidelines at
\url{https://journals.plos.org/plosone/s/materials-and-software-sharing#loc-sharing-code}
and ensure that your code is shared in a way that follows best practice and
facilitates reproducibility and reuse.

\end{cmt}
\vspace{0.1cm}
\noindent
\underline{\textbf{Response:}}
\vspace{0.2cm}

\emph{[TODO]}

\vspace{0.3cm}

%%%%%%%%%%%%% Journal Requirement 3
\begin{cmt}{}{}

We note that the grant information you provided in the `Funding Information'
and `Financial Disclosure' sections do not match.
When you resubmit, please ensure that you provide the correct grant numbers for
the awards you received for your study in the `Funding Information' section.

\end{cmt}
\vspace{0.1cm}
\noindent
\underline{\textbf{Response:}}
\vspace{0.2cm}

\emph{[TODO]}

\vspace{0.3cm}


%%%%%%%%%%%%% Journal Requirement 4
\begin{cmt}{}{}

Thank you for stating the following financial disclosure:

``The author(s) declared that financial support was received for this work
and/or its publication. This research was funded by the Committee of Science of
the Ministry of Science and Higher Education of the Republic of Kazakhstan
(Grant no. BR27198099).''

Please state what role the funders took in the study. If the funders had no
role, please state: ``The funders had no role in study design, data collection
and analysis, decision to publish, or preparation of the manuscript.''

If this statement is not correct you must amend it as needed.

Please include this amended Role of Funder statement in your cover letter; we
will change the online submission form on your behalf.

\end{cmt}
\vspace{0.1cm}
\noindent
\underline{\textbf{Response:}}
\vspace{0.2cm}

\emph{[TODO]}

\vspace{0.3cm}


%%%%%%%%%%%%% Journal Requirement 5
\begin{cmt}{}{}

Please note that your Data Availability Statement is currently missing the
repository name and/or the DOI/accession number of each dataset OR a direct link
to access each database. If your manuscript is accepted for publication, you
will be asked to provide these details on a very short timeline. We therefore
suggest that you provide this information now, though we will not hold up the
peer review process if you are unable.

\end{cmt}
\vspace{0.1cm}
\noindent
\underline{\textbf{Response:}}
\vspace{0.2cm}

\emph{[TODO]}

\vspace{0.3cm}


%%%%%%%%%%%%% Journal Requirement 6
\begin{cmt}{}{}

We notice that your supplementary figures are uploaded with the file type
`Figure'. Please amend the file type to `Supporting Information'. Please ensure
that each Supporting Information file has a legend listed in the manuscript
after the references list.

\end{cmt}
\vspace{0.1cm}
\noindent
\underline{\textbf{Response:}}
\vspace{0.2cm}

\emph{[TODO]}

\vspace{0.3cm}


%%%%%%%%%%%%% Journal Requirement 7
\begin{cmt}{}{}

If the reviewer comments include a recommendation to cite specific previously
published works, please review and evaluate these publications to determine
whether they are relevant and should be cited. There is no requirement to cite
these works unless the editor has indicated otherwise.

\end{cmt}
\vspace{0.1cm}
\noindent
\underline{\textbf{Response:}}
\vspace{0.2cm}

\emph{[TODO]}

\vspace{0.3cm}


%%%%%%%%%%%%%%%%%%%%%%%%%%%%%%%%%%%%%%%%%%%%%%%%%%%%%%%%%%%%%%% Reviewer 1
\newpage
\setcounter{tcb@cnt@cmt}{0}
\section*{Reviewer \#1}

\vspace{0.4cm}

%%%%%%%%%%%% Comment 2.1 -- MR1
\begin{cmt}{}{}

MR1. The binary attention labels have uncertain construct validity. Dataset B
was collected for stress analysis rather than attention measurement, while
Dataset A uses instructed states presented in a fixed temporal sequence.
Drowsiness also involves eye closure. Consequently, the classes confound
attention with stress, arousal, visual input, ocular activity and time-on-task.

\end{cmt}
\vspace{0.1cm}
\noindent
\underline{\textbf{Response:}}
\vspace{0.2cm}

\textbf{[TODO --- this response must be written again from scratch; the previous
draft described a manuscript that does not exist.]} It claimed a renamed Table~1
with operationally defined classes, a paragraph after Table~1 stating what the
labels do not denote, a Limitations subsection treating the Reviewer's five
confounds one by one, a M\"{u}hl et al.\ (2014) framing of stress as an affective
context, and five label-mapping sensitivity arms. None of these are in
\texttt{manuscript.tex}: Table~1 is unchanged (``Unified binary attention label
mapping'', classes ``High Attention'' and ``Low Attention''), there is no
Limitations subsection at all although two passages of Methodology already refer
the reader to one, \verb|muhl2014workload| sits in \texttt{refs.bib} uncited, and
neither the configuration nor the sweep contains any label-mapping arm beyond
\texttt{dataset.exclude\_conditions}.

\textbf{[TODO --- what is true today and may be reused.]} Drowsy epochs are
excluded from the primary analysis, and the exclusion happens at the head of the
preparation pipeline, before channel selection, filtering, ICA and windowing, so
that no drowsy sample reaches the ICA decomposition or the feature scaler. This
holds in the code and is stated in the Data preparation subsection. The two
factual observations about the datasets --- that the authors of Dataset~A report
that no independent objective control could be introduced, and that the public
release of Dataset~B omits the participants' responses --- are also sound and can
be carried over.

\textbf{[TODO --- what the answer has to argue.]} The construct problem is not
solved by renaming the classes. The answer must concede that \emph{attention} is
an umbrella term with no single operational definition, reframe the target
accordingly, and present the label mapping as an explicit operational choice
whose consequences are measured rather than assumed; see
\texttt{thinking\_about\_introduction\_section/} in the repository root. It can
only be written after the corresponding paragraphs exist in the Introduction and
Materials sections of the manuscript.

\vspace{0.3cm}


%%%%%%%%%%%% Comment 2.2 -- MR2
\begin{cmt}{}{}

MR2. The principal within-dataset baseline is affected by substantial
information leakage. One-second windows separated by only 0.125 seconds share
87.5\% of their samples, yet temporally adjacent windows are distributed across
stratified folds. Acknowledging this issue as a limitation does not make the
reported baseline scientifically valid.

\end{cmt}
\vspace{0.1cm}
\noindent
\underline{\textbf{Response:}}
\vspace{0.2cm}

We agree. We reconstructed the within-dataset baseline using group-aware
splitting at the recording-unit level. Thus, all overlapping windows derived
from a physical recording are assigned wholly to either the training or the
test partition of a fold; no temporally adjacent overlapping windows are split
between them. We also replaced the previous one-second windowing scheme with
five-second windows and a 0.5-second overlap. The baseline results are being
recomputed under this leakage-free design.

\textbf{[TODO --- not yet in the manuscript, and not yet recomputed.]} The
recording-unit grouping and the 5\,s window with 0.5\,s overlap exist in the
code, but the Data preparation subsection of the manuscript still describes 1\,s
windows with a 0.125\,s step --- the very geometry the Reviewer objects to --- and
the leakage-free baseline has not been computed: the sweep is still running and
no analysis has been produced from it. Reconcile the manuscript and confirm the
numbers before this response is treated as final.

\vspace{0.3cm}


%%%%%%%%%%%% Comment 2.3 -- MR3
\begin{cmt}{}{}

MR3. Several validation procedures are insufficiently defined. Dataset B
trials appear to be treated as sessions without justification; the
composition of the ``task-only'' datasets is unclear; class counts are
absent; and cross-task testing does not separate task transfer from
participant identity.

\end{cmt}
\vspace{0.1cm}
\noindent
\underline{\textbf{Response:}}
\vspace{0.2cm}

We agree and have corrected the terminology and validation design throughout.

\emph{Dataset B trials treated as sessions.} The SAM40 repeated task
recordings are treated as trials rather than as independent physical
sessions. Cross-session validation is now restricted to the dataset with
documented session entities, whereas SAM40 is evaluated with a separately
named cross-trial protocol.

\emph{Composition of the ``task-only'' datasets.} We now define every
analytical subset explicitly by the conditions retained from the complete
dataset mapping. The exact excluded conditions are stored in the resolved
configuration and in the identity of the prepared data artifact, so that each
task-specific result can be reconstructed unambiguously.

\emph{Missing class counts.} We have added training and test class counts
for every fold, together with per-class support in the classification
tables. The revised Table~2 will report the relevant sample sizes alongside
its performance measures.

\emph{Subject-identity confounding in cross-task testing.} We rebuilt
cross-task validation as leave-one-subject-out transfer: for each fold, the
source-task training data exclude subject $S$, and the target-task test data
contain only subject $S$. Source and target partitions are therefore
subject-disjoint by construction.

\textbf{[TODO --- the class-count part is not true yet.]} Fold composition and
per-class support are computed by the analysis script, but the analysis has not
been run and the manuscript contains no Table~2: the Results section is a cleared
placeholder. The three other parts of this response --- session versus trial
terminology, explicit subset composition, subject-disjoint cross-task transfer
--- are implemented in the code and stated in the manuscript.

\vspace{0.3cm}


%%%%%%%%%%%% Comment 2.4 -- MR4
\begin{cmt}{}{}

MR4. Reproducibility is incomplete within the manuscript. Filter type and order,
ICA implementation and rejection criteria, STFT window and FFT parameters,
feature dimensionality, scaler-fitting scope, XGBoost hyperparameters, tuning
procedure, random seeds and fold construction are not adequately reported.

\end{cmt}
\vspace{0.1cm}
\noindent
\underline{\textbf{Response:}}
\vspace{0.2cm}

We expanded the reproducibility specification. The versioned configuration now
records filtering, ICA method and EOG-proxy criterion, epoching, channel set,
feature extraction, model hyperparameters, random seed, and validation
protocol. Each run retains its resolved configuration. Hyperparameters are
fixed a priori for the benchmark comparison; no data-driven hyperparameter
tuning is performed, and this will be stated explicitly in the Methods.

\textbf{[TODO --- the configuration half is done, the manuscript half is not.]}
The versioned configuration does record filtering, ICA, epoching, channel set,
features, model hyperparameters, seed and protocol, and every run stores its
resolved configuration. The manuscript reports none of it: the Models subsection
is a cleared placeholder, and the model family its text names (XGBoost) still has
to be reconciled with the set of models actually configured --- logistic
regression, XGBoost, EEGNet, EEGConformer and ShallowNet. Write the
Methods first, then finalise this response.

\vspace{0.3cm}


%%%%%%%%%%%% Comment 2.5 -- R1
\begin{cmt}{}{}

R1. The metric reporting is internally inconsistent. In single-label
classification, support-weighted recall is mathematically equal to accuracy.
Table 2 nevertheless reports, for example, 0.81 accuracy and 0.97 weighted
recall for Full B, and 0.73 accuracy and 0.90 weighted recall for its
cross-subject evaluation. Either the averaging description or the results are
incorrect.

\end{cmt}
\vspace{0.1cm}
\noindent
\underline{\textbf{Response:}}
\vspace{0.2cm}

We agree and recompute all reported metrics from the saved window-level
predictions rather than from independently maintained summary values. The
revised analysis exports confusion matrices, per-class precision, recall and
F1, balanced accuracy, macro-F1, MCC, AUROC where defined, and the
training-fold majority-class baseline. The corrected Table~2 will replace the
inconsistent earlier values.

\textbf{[TODO --- no corrected metrics exist yet.]} Recomputation from saved
window-level predictions is implemented in the analysis script, but it has not
been run, and there is no Table~2 in the manuscript to correct. This response
cannot be sent before the recomputed values exist and the inconsistency the
Reviewer identified is demonstrably gone from them.

\vspace{0.3cm}


%%%%%%%%%%%% Comment 2.6 -- R2
\begin{cmt}{}{}

R2. Revise the presentation and interpretation. Figure 1 should include every
validation family with uncertainty; Figure 2 should distinguish diagonal
baselines from transfer cells.

\end{cmt}
\vspace{0.1cm}
\noindent
\underline{\textbf{Response:}}
\vspace{0.2cm}

We have redesigned Figure~1 to report every validation family and to show
participant-level 95\% confidence intervals for the displayed metric. The
figure is generated directly from the recomputed results, rather than from a
manually maintained plotting table.

We have also redesigned Figure~2 as a source--target transfer matrix. Matched
within-composition baselines occupy the diagonal and are visually marked as
such; all off-diagonal cells represent transfer. This separates the relevant
within-condition reference from cross-condition performance.

\textbf{[TODO --- both figures are code, not manuscript.]} The two designs
described here are implemented in the analysis script, but neither figure has
been generated, and the Results section of the manuscript currently contains no
figures at all. The past tense is justified only once they are produced and
included.

\vspace{0.3cm}


%%%%%%%%%%%% Comment 2.7 -- R3
\begin{cmt}{}{}

R3. The 15 references are relevant and include the original publications for
both datasets, established cognitive-task literature and selected studies on EEG
transfer. However, the coverage is too limited for the broad methodological
claims advanced. Statements concerning DANN, ADDA, CORAL, TCA, Riemannian
alignment, transformer architectures, self-supervised learning, XGBoost
performance and low-density ICA are presented without appropriate citations.

\end{cmt}
\vspace{0.1cm}
\noindent
\underline{\textbf{Response:}}
\vspace{0.2cm}

We thank the reviewer for this comment, which prompted a full audit of the
bibliography rather than a simple expansion of it. We have acted on it in two
ways.

\textbf{(a) Expanded framing and citations in the Introduction.} The
Introduction previously moved directly from the general difficulty of EEG
generalization to the aims of the study, without grounding the central
hypothesis. We have added two paragraphs. The first states why cross-paradigm
transfer is expected to be fragile: transfer learning depends on domain
\emph{and label} compatibility between source and target, and attention lacks a
single accepted operational definition, so datasets nominally labelled for
``attention'' may in fact probe different attentional types, producing label
mismatch under domain shift. The second states the gap explicitly, adopts the
term \emph{zero-shot generalization} for what is measured here, and positions
the study as a baseline against which transfer-learning and domain-adaptation
methods can be compared. These paragraphs add five references that were absent
from the submitted version (Wu et al. 2022; Zhang et al. 2020; Hommel et al.
2019; Sun et al. 2025; Li et al. 2025).

\textbf{(b) Verification of the entire reference list.} Prompted by the
reviewer's remark on reference quality, every entry in the bibliography was
resolved against Crossref, arXiv or the publisher record. This uncovered errors
in the submitted reference list that we correct here and regret. Specifically:
the cross-task workload study cited in Background was printed with an author
list that does not belong to it (the correct authors are Ke et al., Front
Neurosci 2021;15:703139); the single-channel Kalman/ELM study was printed with
an incorrect author list, journal, year and DOI (the correct record is He et
al., Front Hum Neurosci 2025;18:1481493); and three further cited entries could
not be resolved to any existing publication and have been replaced by verified
sources supporting the same statements (Esterman et al., Cereb Cortex
2013;23:2712--2723 for sustained attention; Kitaura et al., Clin Neurophysiol
Pract 2017;2:193--200 for attentional demand during mental arithmetic; Tomasino
and Gremese, Front Hum Neurosci 2016;9:693 for the mental-rotation network).
Duplicate entries for the same work, and several entries with incorrect author
names, volumes or DOIs, were also corrected. The revised reference list has been
checked to contain no unresolved or duplicated records.

We note that part of this comment remains open. The passages discussing
domain-adversarial and adversarial discriminative adaptation, correlation
alignment, transfer component analysis, Riemannian alignment, transformer-based
and self-supervised representation learning, gradient-boosting performance on
EEG features, and the stability of ICA on low-density montages are still stated
without supporting citations. We will supply appropriate references for each of
these claims in the next revision round, together with the specific recent works
requested by Reviewer~3.

\textbf{[TODO --- two uncited entries, and Background does not exist.]} Part (a)
is accurate: the two Introduction paragraphs and their five citations are in the
manuscript. \verb|verma2023attentionet| and \verb|muhl2014workload|, however, are
present in \texttt{refs.bib} but cited nowhere, so they do not enter the
reference list; either cite them or drop them, and keep the count above in
agreement with what is actually cited. The Background section --- where the
missing citations for domain adaptation, alignment, transformers,
self-supervised learning, gradient boosting and low-density ICA belong --- is a
cleared placeholder, which is why the closing paragraph of this response is
accurate today; rewrite it once that section exists.

\vspace{0.3cm}


%%%%%%%%%%%%%%%%%%%%%%%%%%%%%%%%%%%%%%%%%%%%%%%%%%%%%%%%%%%%%%% Reviewer 2
\newpage
\setcounter{tcb@cnt@cmt}{0}
\section*{Reviewer \#2}

\vspace{0.4cm}

%%%%%%%%%%%% Comment 3.1
\begin{cmt}{}{}

Reconsider the target construct. The problem should be presented either as
discrimination between operationally defined experimental conditions or as
attention classification supported by independent attention measurements.

\end{cmt}
\vspace{0.1cm}
\noindent
\underline{\textbf{Response:}}
\vspace{0.2cm}

We have adopted the first of the two formulations the Reviewer permits, and the
revised manuscript is built around it throughout.

The second formulation is unavailable for these data, and the manuscript now
demonstrates this rather than assuming it. The authors of Dataset~A report that
no independent objective control could be introduced, because active probing of
a disengaged participant destroys the state being measured and no external
correlate of passive focus could be identified. The public release of Dataset~B
comprises raw and filtered recordings, artifact-removed signals, electrode
coordinates and per-trial subjective stress ratings, but not the participants'
responses, so performance-based labels cannot be constructed from it. Neither
dataset can therefore support attention classification anchored to an
independent measurement of attention.

\textbf{[TODO --- the manuscript does not yet carry this formulation.]} What
exists: the Methods state that the model was trained to predict, from the
spectral features of a window, whether the participant was completing a
high-demand or a low-demand condition during that window, and not to predict a
level of attention. What does not exist: Table~1 still carries the classes
``High Attention'' and ``Low Attention'', no paragraph follows it stating what
the labels do not denote, and the Abstract is a cleared placeholder, so it says
nothing about the sense in which ``attention decoding'' is used in the title.
The five label-mapping arms have no counterpart in the code either --- see
Comment~3.2. Rewrite this response together with the response to MR1, once the
corresponding manuscript paragraphs are written.

We would add that this constraint is not peculiar to our two datasets. The same
criticism has been published against vigilance and workload decoders whose
labels encode the condition of the protocol rather than the state of the
participant. The revision makes this the subject of the study rather than a
caveat to it: since publicly available attention datasets supply protocol-defined
labels, the question that can honestly be answered is how much of a decoder's
performance survives when the operational definition of its labels is changed.

\vspace{0.3cm}


%%%%%%%%%%%% Comment 3.2
\begin{cmt}{}{}

Sensitivity analyses should exclude drowsy or eyes-closed epochs and evaluate
alternative label mappings.

\end{cmt}
\vspace{0.1cm}
\noindent
\underline{\textbf{Response:}}
\vspace{0.2cm}

We agree that drowsiness should not be silently merged with awake unfocused
behaviour. The primary analysis now excludes the explicitly labelled drowsy
segment, and alternative condition mappings are represented as separate,
reproducible analyses. We cannot validly identify or exclude eyes-closed epochs
from the available recordings: they contain no EOG channel, video, or
epoch-level eye-state annotation. Frontal EEG channels were used only as proxies
to identify EOG-like ICA components; they do not provide ground truth for eye
closure. We therefore do not claim eyes-closed epoch rejection, because an
unvalidated amplitude-based proxy could remove meaningful EEG activity.

\textbf{[TODO --- the alternative mappings do not exist.]} There is no
label-mapping sensitivity analysis anywhere: a run selects its conditions through
\texttt{dataset.exclude\_conditions}, and the sweep defined in the Makefile varies
the validation protocol, the dataset and the model --- never the mapping. No
M0--M4 arms are defined in the manuscript or in the configuration, and none are
being computed. Either implement the arms and describe them in the manuscript, or
answer this comment with the narrower thing that is true --- the drowsy exclusion
above, and the per-task subsets already used for cross-task transfer --- and say
plainly that a full mapping-sensitivity analysis is outside the scope of this
revision.

\vspace{0.3cm}


%%%%%%%%%%%% Comment 3.3
\begin{cmt}{}{}

Reconstruct all validation splits at the participant, session, or trial level
before window segmentation.

\end{cmt}
\vspace{0.1cm}
\noindent
\underline{\textbf{Response:}}
\vspace{0.2cm}

We reconstructed all validation protocols around recording-level metadata
(participant, documented session, or trial) before assigning windows to folds.
Every window from a recording unit remains on one side of a split. The revised
Methods distinguish cross-session and cross-trial validation rather than using
the terms interchangeably.

\vspace{0.3cm}


%%%%%%%%%%%% Comment 3.4
\begin{cmt}{}{}

Every preprocessing transformation, including scaling and data-dependent feature
processing, should be fit exclusively using the training data.

\end{cmt}
\vspace{0.1cm}
\noindent
\underline{\textbf{Response:}}
\vspace{0.2cm}

All learned transformations, including standardization, are contained in the
estimator pipeline and are fit separately within each training fold. The
handcrafted feature transforms are deterministic functions of an individual
window and estimate no parameter from the data; they therefore have no fitting
scope that could transmit information across folds.

\vspace{0.3cm}


%%%%%%%%%%%% Comment 3.5
\begin{cmt}{}{}

Cross-task experiments should also use subject-disjoint source and target
partitions.

\end{cmt}
\vspace{0.1cm}
\noindent
\underline{\textbf{Response:}}
\vspace{0.2cm}

We agree and rebuilt cross-task validation as leave-one-subject-out transfer.
For each held-out participant, the source-task training partition excludes that
participant and the target-task test partition contains only that participant.
The source and target partitions are consequently subject-disjoint.

\vspace{0.3cm}


%%%%%%%%%%%% Comment 3.6
\begin{cmt}{}{}

Audit and recompute all metrics from the saved predictions. Class counts,
confusion matrices, per-class precision and recall, balanced accuracy, macro-F1,
MCC, and AUROC should be reported alongside majority-class baselines.

\end{cmt}
\vspace{0.1cm}
\noindent
\underline{\textbf{Response:}}
\vspace{0.2cm}

All performance measures are now recomputed from saved window-level
predictions. The revised outputs include fold composition, confusion matrices,
per-class precision, recall and F1, balanced accuracy, macro-F1, MCC, AUROC
where defined, and majority-class baselines. Table~2 will report these metrics
and their definitions consistently.

\textbf{[TODO --- implemented, but neither computed nor reported.]} Every
measure listed here is implemented in the analysis script and none of it is
invented; none has been produced yet, because the sweep is still running and the
analysis has not been run over it. The manuscript has no Table~2 at present.

\vspace{0.3cm}


%%%%%%%%%%%% Comment 3.7
\begin{cmt}{}{}

Fold- or participant-level distributions, confidence intervals, and appropriate
paired statistical comparisons are also required.

\end{cmt}
\vspace{0.1cm}
\noindent
\underline{\textbf{Response:}}
\vspace{0.2cm}

The revised analysis reports fold-level and participant-level distributions and
95\% confidence intervals with their aggregation unit stated explicitly. It
also performs paired, participant-level Wilcoxon comparisons only when runs
share the same target data, with Holm and false-discovery-rate adjustments for
multiple comparisons.

\textbf{[TODO --- as for Comment~3.6.]} Fold- and participant-level
distributions, 95\% confidence intervals and paired Wilcoxon tests with Holm and
false-discovery-rate correction exist in the analysis script; they have not been
run, and the manuscript reports none of them.

\vspace{0.3cm}


%%%%%%%%%%%% Comment 3.8
\begin{cmt}{}{}

Provide a complete and reproducible description of the methods, including filter
characteristics, the ICA procedure, STFT settings, band-power calculations,
feature dimensions, hyperparameters, tuning strategy, and random seeds.

\end{cmt}
\vspace{0.1cm}
\noindent
\underline{\textbf{Response:}}
\vspace{0.2cm}

We have expanded the Methods and the versioned configuration to specify the
filtering recipe, ICA method and EOG-proxy criterion, window geometry, channel
selection, feature extraction, model hyperparameters, random seed, and fold
construction. Hyperparameters are fixed before the benchmark comparison; no
data-driven tuning is performed. Each run stores its resolved configuration so
that the exact analysis can be reproduced.

\textbf{[TODO --- see the response to MR4.]} The configuration side is done, the
Methods side is not: the Models subsection of the manuscript is a cleared
placeholder, and the model family named in its text contradicts the models
actually configured.

\vspace{0.3cm}


%%%%%%%%%%%% Comment 3.9
\begin{cmt}{}{}

At least simple linear and tree-based benchmark models should be included.

\end{cmt}
\vspace{0.1cm}
\noindent
\underline{\textbf{Response:}}
\vspace{0.2cm}

We added logistic regression as a linear benchmark and gradient-boosted decision
trees (XGBoost) as a tree-based benchmark, evaluated under the same validation
protocols as the deep models.

\textbf{[TODO --- confirm the final benchmark set, then report it.]} The models
configured today are logistic regression, XGBoost, EEGNet, EEGConformer and
ShallowNet; the SVM and ExtraTrees named in earlier drafts of this
letter are no longer in \texttt{configs/model/}. The Models subsection of the
manuscript names none of them, and no results exist for them yet.

\vspace{0.3cm}


%%%%%%%%%%%% Comment 3.10
\begin{cmt}{}{}

An alignment or domain-adaptation baseline would further strengthen the transfer
analysis.

\end{cmt}
\vspace{0.1cm}
\noindent
\underline{\textbf{Response:}}
\vspace{0.2cm}

\emph{[TODO]}

\vspace{0.3cm}


%%%%%%%%%%%% Comment 3.11
\begin{cmt}{}{}

Revise the presentation and interpretation. Figure 1 should include every
validation category with corresponding uncertainty estimates.

\end{cmt}
\vspace{0.1cm}
\noindent
\underline{\textbf{Response:}}
\vspace{0.2cm}

Figure~1 has been redesigned to include every validation family and
participant-level 95\% confidence intervals, generated directly from the
recomputed results.

\textbf{[TODO --- see the response to Comment~R2 of Reviewer~1.]} The figure is
implemented in the analysis script, has not been generated, and is not in the
manuscript.

\vspace{0.3cm}


%%%%%%%%%%%% Comment 3.12
\begin{cmt}{}{}

Figure 2 should distinguish diagonal baseline results from transfer cells.

\end{cmt}
\vspace{0.1cm}
\noindent
\underline{\textbf{Response:}}
\vspace{0.2cm}

Figure~2 has been redesigned as a source--target matrix in which matched
within-composition baselines occupy and are explicitly marked on the diagonal;
off-diagonal cells represent transfer.

\textbf{[TODO --- see the response to Comment~R2 of Reviewer~1.]} The figure is
implemented in the analysis script, has not been generated, and is not in the
manuscript.

\vspace{0.3cm}


%%%%%%%%%%%% Comment 3.13
\begin{cmt}{}{}

Table 2 should include sample sizes and corrected metric definitions.

\end{cmt}
\vspace{0.1cm}
\noindent
\underline{\textbf{Response:}}
\vspace{0.2cm}

The revised Table~2 will include the number of participants, windows, and
class support for each analysis, together with metrics recomputed from the
saved predictions and clearly defined aggregation units.

\textbf{[TODO --- there is no Table~2 yet.]} The Results section of the
manuscript is a cleared placeholder; the sample sizes and metric definitions
promised here can be written only after the analysis has been run.

\vspace{0.3cm}


%%%%%%%%%%%% Comment 3.14
\begin{cmt}{}{}

Claims regarding hardware, causal language, references, and the secondary
data-use ethics statement should also be corrected.

\end{cmt}
\vspace{0.1cm}
\noindent
\underline{\textbf{Response:}}
\vspace{0.2cm}

\emph{[TODO]}

\vspace{0.3cm}


%%%%%%%%%%%%%%%%%%%%%%%%%%%%%%%%%%%%%%%%%%%%%%%%%%%%%%%%%%%%%%% Reviewer 3
\newpage
\setcounter{tcb@cnt@cmt}{0}
\section*{Reviewer \#3}

\vspace{0.4cm}

%%%%%%%%%%%% Comment 4.1
\begin{cmt}{}{}

I have carefully reviewed the article. At a first glance, this paper would have
been really very interesting but unfortunately it is really too short and
superficial for the promises it makes at the beginning. The authors set up a
compelling premise in the introduction, promising to systematically evaluate how
neural markers of attention transfer across diverse paradigms and validation
strategies. However, the execution falls completely flat. It would be necessary
to compare all the most advanced deep models for EEG decoding at the state of
the art and see if, by passing from one task to another, they do better or
worse. Instead, here unfortunately there are very few discussions on state of
the art models and on datasets used for many tasks such as motor imagery, sleep,
and other cognitive states. The entire methodology relies on a simplistic
feature extraction pipeline combined with a traditional machine learning
classifier, which completely ignores the massive advancements made in deep
representation learning for electroencephalography over the last several years.

\end{cmt}
\vspace{0.1cm}
\noindent
\underline{\textbf{Response:}}
\vspace{0.2cm}

We appreciate this concern. The revised benchmark evaluates classical
feature-based models alongside EEG-specific deep architectures --- EEGNet,
EEGConformer and ShallowNet --- which receive the prepared EEG windows
directly rather than handcrafted spectral features and are evaluated under the
same validation protocols. We have also narrowed the manuscript's claims to the
two available datasets and their operationally defined conditions.

\textbf{[TODO --- the results and the state-of-the-art discussion are both
missing.]} The deep models are configured and are being run, but no results exist
for them, and the Background section that would review deep EEG decoding is a
cleared placeholder. This response cannot stand as written until both are in
place.

\vspace{0.3cm}


%%%%%%%%%%%% Comment 4.2
\begin{cmt}{}{}

When discussing the state of the art in electroencephalography decoding, it is
entirely unacceptable to limit the analysis to gradient boosted decision trees
like XGBoost fed with basic short time Fourier transform spectral power
features. The field has moved well beyond these classical baselines. Advanced
deep models such as convolutional neural networks specifically designed for time
series, spatial temporal graph convolutional networks, and transformer based
architectures have become the standard for brain computer interfaces and
cognitive state monitoring. The manuscript barely touches upon these
architectures, briefly mentioning transformers in the related work but failing
to implement or benchmark any of them. To truly understand whether neural
signatures of attention are invariant across paradigms, one must leverage models
capable of learning hierarchical spatial and temporal representations directly
from the raw or minimally preprocessed signals. By relying on predefined
spectral bands and overlapping windows, the authors impose a rigid inductive
bias that likely destroys the subtle, task invariant dynamics they claim to be
looking for. Therefore, the conclusion that current classifiers capture context
dependent features rather than invariant signatures is heavily confounded by the
fact that the authors used a very rudimentary classifier that lacks the capacity
to learn invariant representations in the first place.

\end{cmt}
\vspace{0.1cm}
\noindent
\underline{\textbf{Response:}}
\vspace{0.2cm}

We added EEG-specific deep architectures --- EEGNet, EEGConformer and
ShallowNet --- as end-to-end benchmarks on the prepared raw EEG windows, alongside the
feature-based logistic regression and XGBoost baselines. This comparison tests
whether the observed validation-pattern differences persist when models can learn
representations directly from the signal rather than only from predefined
spectral features.

\textbf{[TODO --- as for Comment~4.1, no results exist yet.]} Earlier drafts of
this letter also named SVM and ExtraTrees, which are no longer configured; keep
the list here in agreement with \texttt{configs/model/} and with the Models
subsection of the manuscript once that subsection is written.

\vspace{0.3cm}


%%%%%%%%%%%% Comment 4.3
\begin{cmt}{}{}

Furthermore, the narrow focus on a passive control task and short cognitive
tests like the Stroop test and mental arithmetic severely limits the scope of
the claims regarding cross task generalization. A comprehensive evaluation
should have included datasets spanning a much wider variety of neurological
tasks. For instance, motor imagery is one of the most widely studied paradigms
in the field, and analyzing how attentional engagement during motor imagery
transfers to active cognitive tests or sleep staging would have provided a much
stronger foundation for the paper. Sleep datasets, in particular, offer a
continuous spectrum of arousal and attentional states that could have perfectly
complemented the drowsy condition in the first dataset. The lack of discussion
on these broader applications and the datasets commonly used for them makes the
paper feel isolated from the broader brain computer interface community and
leaves the reader wondering how these models would perform in real world
clinical scenarios.

\end{cmt}
\vspace{0.1cm}
\noindent
\underline{\textbf{Response:}}
\vspace{0.2cm}

We do not extend the evaluation to motor imagery or sleep staging, and we want
to be explicit about why this is a scope boundary rather than a limitation we
intend to resolve with more data. Motor imagery and sleep staging are defined
by different operational constructs than the sustained-attention and
cognitive-load states this paper targets: motor imagery decoding identifies
sensorimotor-rhythm modulation tied to imagined movement, and sleep staging
assigns polysomnography-defined vigilance stages. Neither is a measure of
attentional engagement during an ongoing task. Adding them would not test
whether attention-decoding models transfer across attention-related paradigms;
it would test transfer across constructs the models were never designed to
represent, which is a different research question from the one this paper asks.
We therefore keep cross-task and cross-dataset evaluation within operationally
defined attention-related paradigms and state this as an intentional scope
boundary in the manuscript, rather than implying that a wider survey of
unrelated BCI paradigms would have strengthened the present claims.

\textbf{[TODO --- this scope boundary is nowhere in the manuscript.]} The
argument stands on its own and needs no new data, but no section currently makes
it: the Background section, where it belongs together with the discussion of the
motor-imagery and sleep paradigms the Reviewer asks about, is a cleared
placeholder.

\vspace{0.3cm}


%%%%%%%%%%%% Comment 4.4
\begin{cmt}{}{}

To rectify these significant omissions, the authors must update their literature
review and theoretical framework by citing and discussing recent works that
successfully apply deep learning to complex electroencephalography tasks. I
strongly require the authors to cite the paper titled \textit{Investigating the
impact of rational dilated wavelet transform on motor imagery EEG decoding with
deep learning models} published in the year 2025. It is extremely important to
cite this work because it clearly demonstrates how advanced time frequency
transforms can be integrated with deep learning models to significantly improve
decoding performance on complex tasks like motor imagery. This highlights the
inadequacy of the standard short time Fourier transform approach used in the
current manuscript. Additionally, the authors must cite the paper titled
\textit{RatioWaveNet: A Learnable RDWT Front-End for Robust and Interpretable
EEG Motor-Imagery Classification}, also published in 2025. Citing this paper is
crucial because it showcases the current state of the art in developing robust,
interpretable front ends that are completely learnable, effectively replacing
the static feature extraction pipelines that the authors currently rely on.
Since these papers are very recent, they represent the cutting edge of the field
and provide exactly the kind of state of the art methodological context that
this superficial manuscript currently lacks.

\end{cmt}
\vspace{0.1cm}
\noindent
\underline{\textbf{Response:}}
\vspace{0.2cm}

\emph{[TODO]}

\vspace{0.3cm}


%%%%%%%%%%%% Comment 4.5
\begin{cmt}{}{}

Moving to the semantic and conceptual issues within the manuscript, there are
several major flaws that undermine the validity of the unified labeling scheme.
The authors decided to merge the drowsy state, the unfocused state, and the
resting state into a single low attention class. This is a severe semantic
error. From a neurophysiological perspective, drowsiness involves a transition
into early stage sleep, characterized by the emergence of theta waves, vertex
sharp waves, and a general withdrawal of environmental awareness. This is
fundamentally different from an awake but unfocused state, which might involve
default mode network activation and mind wandering, or a resting state with eyes
open or closed, which is typically dominated by posterior alpha oscillations. By
grouping these distinct physiological phenomena into a single binary class, the
authors create an artificially heterogeneous category that confuses the
classifier. This semantic misalignment invalidates many of the cross dataset
conclusions, as the model is likely learning to distinguish arousal levels
rather than attention levels.

\end{cmt}
\vspace{0.1cm}
\noindent
\underline{\textbf{Response:}}
\vspace{0.2cm}

We accept this criticism and have acted on it rather than argued against it. We
agree that drowsiness is not interchangeable with an awake unfocused state or
with eyes-open rest, and that the neurophysiological grounds the Reviewer gives
--- the transition towards early-stage sleep, with its own signature and with
sustained eye closure --- are sufficient reason not to place these conditions in
one class.

Drowsy epochs are therefore excluded from the primary analysis entirely, and not
reassigned or flagged as a limitation. The exclusion is implemented as a step of
the data-preparation pipeline, ahead of channel selection, filtering, artifact
removal and windowing, so that drowsy data do not enter the ICA decomposition or
the fitting of the feature scaler either. The low-demand class of the primary
mapping now contains only the unfocused and resting conditions, and the
manuscript states the Reviewer's own reasoning as the justification for the
exclusion.

\textbf{[TODO --- the remainder of this response describes a manuscript that does
not exist.]} Table~1 has not been renamed: it is still ``Unified binary attention
label mapping'' with the classes ``High Attention'' and ``Low Attention'', and no
paragraph follows it. There is no sensitivity arm~M0, in the manuscript or in the
code, so the Reviewer's inference about arousal is still not testable. There is
no Limitations subsection at all, although two passages of Methodology already
refer the reader to one. Write the Table~1 reframe and the Limitations subsection
first; only then can this response state what the arousal objection is actually
answered with.

\vspace{0.3cm}


%%%%%%%%%%%% Comment 4.6
\begin{cmt}{}{}

Furthermore, framing mental arithmetic purely as an attention task ignores the
massive working memory load it requires, which produces entirely different
cortical activation patterns compared to the simple visual monitoring required
in the train simulator task.

\end{cmt}
\vspace{0.1cm}
\noindent
\underline{\textbf{Response:}}
\vspace{0.2cm}

We accept the observation and have not attempted to argue it away. Under the
operational definition adopted in this revision, the working-memory load of
mental arithmetic is not a confound to be removed from an attention contrast but
a constituent of the high-demand condition, and mental arithmetic is no longer
framed as a pure attention task.

We also draw the broader consequence that follows from the Reviewer's point. The
attention-decoding and mental-workload-decoding literatures operationalize their
targets in the same way --- task versus rest, or high versus low demand --- and
are to that extent operationally indistinguishable. This is a further instance of
the main argument of the revision rather than an exception to it: where the label
is supplied by the protocol condition, the boundary between the two constructs is
not drawn by the data.

\textbf{[TODO --- three claims were removed from this response, each
unsupported.]} (i) That cross-task transfers involving arithmetic were
consistently the weakest: those numbers come from the superseded pipeline, were
removed from the manuscript, and have not been recomputed. (ii) Label-mapping
arm~M4, which would keep the Dataset~B tasks separate: no such arm exists in the
configuration or in the sweep, and nothing of the kind is being computed. (iii)
The alpha-desynchronization and frontal-midline-theta rationale for aligning the
two protocols: the manuscript still argues that alignment from the psychological
labels of the individual tasks, which is exactly what the Reviewers did not
accept. Rewrite this response once the manuscript carries a mechanism-based
rationale and once cross-task numbers exist.

\vspace{0.3cm}


%%%%%%%%%%%% Comment 4.7
\begin{cmt}{}{}

Similarly, the grammatical and syntactic quality of the manuscript leaves much
to be desired. There are numerous instances of poor phrasing and missing
punctuation that disrupt the reading flow and obscure the scientific meaning.
For example, in the description of Dataset A, the authors write that all
experiments were performed with volunteer participants chosen among students.
This phrasing is awkward and unscientific; it would be much better to state that
the participant cohort consisted of university student volunteers. In the
description of Dataset B, there are glaring punctuation errors in the
descriptions of the tasks. The text reads Stroop Color Word Test participants
were required to identify the color, and Arithmetic problem solving participants
mentally solved a mathematical problem, and Mirror Image Recognition
participants were shown images. In all these cases, there is a missing colon or
em dash after the name of the task, making it seem like the participants
themselves are named after the tasks. These are careless grammatical errors that
should have been caught during proofreading. Furthermore, the phrasing in the
results section, specifically the sentence stating that accuracy and F1 score
reaching 0.86 for Dataset A and 0.81 and 0.88 respectively for Dataset B, is
disjointed and confusingly structured. The authors also failed to properly
format the cross dataset directions in their text, sometimes omitting the
necessary words to explain the direction of the transfer.

\end{cmt}
\vspace{0.1cm}
\noindent
\underline{\textbf{Response:}}
\vspace{0.2cm}

\emph{[TODO]}

\vspace{0.3cm}


%%%%%%%%%%%% Comment 4.8
\begin{cmt}{}{}

Another critical methodological failure that must be addressed is the cross task
validation experimental design. The authors openly admit in the limitations
section that in their cross task experiments, the source and target tasks were
drawn from the same pool of Dataset B participants, and consequently, task
transfer and subject identity are not fully disentangled. This is not just a
minor limitation; it is a fatal flaw for a study whose entire premise is about
evaluating transferability and generalization. If the model is evaluated on a
new task but with the same subjects it was trained on, any successful transfer
could merely be the result of the classifier learning subject specific noise
profiles, baseline anatomical features, or electrode impedance artifacts, rather
than genuine task invariant attention markers. A true cross task evaluation must
be performed using a leave one subject out approach across tasks, ensuring that
the model is tested on both unseen tasks and unseen subjects simultaneously, or
at the very least, evaluated on completely disjoint subsets of subjects.

\end{cmt}
\vspace{0.1cm}
\noindent
\underline{\textbf{Response:}}
\vspace{0.2cm}

We agree and rebuilt the cross-task protocol as leave-one-subject-out transfer.
In each fold, source-task training data exclude the held-out participant and
target-task testing uses only that participant. The revised design therefore
tests both an unseen task composition and an unseen participant.

\vspace{0.3cm}


%%%%%%%%%%%% Comment 4.9
\begin{cmt}{}{}

The authors also acknowledge that their within dataset baselines are likely
overly optimistic due to the use of densely overlapping windows. They state that
overlapping one second windows with a zero point one two five second step size
lead to highly autocorrelated samples, and that splitting these windows randomly
across cross validation folds causes temporal leakage. While it is good that
they acknowledge this, merely stating it in the limitations is insufficient for
a paper focused on validation strategies. The core contribution of this paper is
supposed to be a rigorous evaluation of validation schemes, yet they purposefully
include a flawed baseline that they know is leaking data. They should have
implemented the group aware splitting by trial or session that they suggest as
future work directly in this manuscript. Presenting inflated within dataset
metrics undermines the entire comparative analysis, as the drop in performance
observed in cross subject and cross session scenarios might just be the result
of removing data leakage rather than actual domain shift. The failure to correct
this known issue makes the baseline metrics virtually meaningless.

\end{cmt}
\vspace{0.1cm}
\noindent
\underline{\textbf{Response:}}
\vspace{0.2cm}

We agree and removed this leakage from the baseline. Overlapping windows are
grouped by their recording unit, and group-aware cross-validation assigns every
window from a unit to either training or testing, never both. The revised
baseline results are recomputed with this procedure.

\textbf{[TODO --- the manuscript still describes the old windows.]} Group-aware
splitting is implemented, but the Data preparation subsection still states 1\,s
windows with a 0.125\,s step, which is the geometry the Reviewer objects to; the
configuration now uses 5\,s windows with 0.5\,s overlap. The recomputed baseline
does not exist yet.

\vspace{0.3cm}


%%%%%%%%%%%% Comment 4.10
\begin{cmt}{}{}

Additionally, the reliance on a reduced twelve channel montage significantly
limits the potential of the study. While aligning the spatial resolution of the
two datasets is a logical step for cross dataset evaluation, dropping from
thirty two channels down to twelve discards a massive amount of spatial
information that could be vital for deep learning models to learn invariant
representations. The authors note that this low density configuration limits the
stability of the independent component analysis used for artifact removal, which
introduces yet another confounding variable. If the artifacts were not cleanly
removed due to the lack of spatial channels, the classifier might simply be
tracking eye movements, blinks, or muscle tension rather than cortical attention
networks. This is especially problematic given that the Stroop test and mental
arithmetic often elicit different physical responses, facial microexpressions,
and stress induced muscle artifacts compared to a passive train driving
simulator. The failure to adequately separate these physiological artifacts from
true neural signals casts doubt on whether the model is actually decoding
attention at all.

\end{cmt}
\vspace{0.1cm}
\noindent
\underline{\textbf{Response:}}
\vspace{0.2cm}

We clarify that the shared 12-channel montage is selected only after artifact
correction. For each recording unit, ICA is fitted to, and applied on, all EEG
channels available in that original recording; channel selection occurs only
after ICA application and before epoching and model fitting. The common-channel
restriction therefore does not reduce the dimensionality of the ICA
decomposition. It is used only to give the cross-dataset models an identical
set of input channels.

We nevertheless acknowledge an important residual limitation. Neither source
dataset provides an independent EOG channel, so ocular-component candidates are
identified using pre-specified frontal EEG channels as EOG proxies. This permits
consistent automated component selection, but it is not ground-truth validation
of ocular-artifact removal. We have made this limitation explicit and have
restricted our interpretation accordingly: the analyses compare classification
of operationally defined experimental conditions and do not establish that the
models isolate cortical attention signals from all possible physiological
artifacts.

\textbf{[TODO --- the manuscript contradicts this response.]} The code does what
is described here: ICA is fitted and applied on all channels of the original
recording, and the common 12-channel montage is picked afterwards
(\texttt{src/preparation.py}). The Data preparation subsection of the manuscript
states the opposite order --- channel selection first, ICA afterwards. Fix the
manuscript before this response is sent.

\vspace{0.3cm}


%%%%%%%%%%%% Comment 4.11
\begin{cmt}{}{}

Finally, the absence of any discussion regarding model explainability or feature
importance is a missed opportunity. Even with a simpler model like XGBoost, the
authors could have analyzed which spectral bands and channels were most critical
for the cross task and cross dataset predictions. Without an analysis of the
underlying features driving the decisions, the claim that models capture context
dependent features remains purely speculative. State of the art deep models,
such as the ones recommended for citation earlier, often include interpretability
modules or attention mechanisms that allow researchers to visualize the spatial
and temporal features the network relies on. By neglecting this aspect, the
authors fail to provide any neuroscientific insights into why the transfers
failed, leaving the paper as a purely empirical exercise in applied machine
learning with poor experimental controls.

\end{cmt}
\vspace{0.1cm}
\noindent
\underline{\textbf{Response:}}
\vspace{0.2cm}

We added fold-wise feature attribution for the models that expose it directly:
coefficients for logistic regression and gain-based importances for XGBoost. The
analysis maps these values back to the configured channel order and frequency
bands and reports their fold average. We present this as model-specific
descriptive evidence, not as an explanation of the deep models.

\textbf{[TODO --- saved, but not yet aggregated or reported.]} Per-fold
attributions are stored with every run, but the analysis that aggregates them has
not been run and the manuscript has no section reporting them. Earlier drafts of
this letter named ExtraTrees, which is no longer configured.

\vspace{0.3cm}


%%%%%%%%%%%% Comment 4.12
\begin{cmt}{}{}

In conclusion, the manuscript fails to deliver on its ambitious initial
premises. To be considered for publication, it requires a complete overhaul. The
authors must replace their obsolete feature extraction and gradient boosting
pipeline with a suite of advanced deep learning architectures, comparing their
performance across various transfer scenarios. They must expand their dataset
inclusion to encompass a wider variety of cognitive states and tasks,
particularly motor imagery and sleep stages, to provide a true measure of
paradigm invariance. They must correct the severe semantic errors in their label
harmonization process, particularly the grouping of sleep related states with
resting states, and fix all grammatical mistakes and punctuation omissions
throughout the text. Most importantly, they must implement rigorous, leakage
free validation schemes for all their baselines and ensure that cross task
evaluations are completely separated from subject identity confounds. Without
these major revisions and a significant expansion in scope, the paper remains
too superficial and methodologically flawed to contribute meaningfully to the
state of the art in electroencephalography decoding.

\end{cmt}
\vspace{0.1cm}
\noindent
\underline{\textbf{Response:}}
\vspace{0.2cm}

We appreciate the detailed assessment. The revision replaces the leaking
baseline and the subject-confounded cross-task protocol, recomputes metrics from
saved predictions, adds participant-level uncertainty and paired comparisons, and
benchmarks linear and tree-based models against EEG-specific deep architectures
(EEGNet, EEGConformer, ShallowNet). We also exclude the explicitly
labelled drowsy segment from the primary analysis and narrow the interpretation
to operationally defined attention-related conditions, deliberately outside the
scope of unrelated BCI constructs such as motor imagery or sleep staging. The
remaining detailed responses specify the corresponding changes and limitations.

\textbf{[TODO --- summary response, to be finalised last.]} It can only be
written once the individual responses above are settled: the sweep is still
running, no analysis has been produced from it, and the Abstract, Background,
Models and Results sections of the manuscript are cleared placeholders.

\vspace{0.3cm}

\end{document}
